\documentclass[12pt]{article}
\usepackage[utf8]{inputenc}
\usepackage[T1]{fontenc}
\usepackage[margin=1in]{geometry}
\usepackage{setspace}
\usepackage{xcolor}
\usepackage[colorlinks=true, urlcolor=blue!60!black, linkcolor=blue!60!black]{hyperref}
\setlength{\parindent}{2em}
\setlength{\parskip}{1.25em}
\renewcommand{\baselinestretch}{1.0}
\begin{document}
\begin{center}
{\LARGE \textit{Miller v. California}, 413 U.S. 15 (1973)\par}
\vspace{0.8em}
{\large Case Brief\par}
\end{center}
\vspace{1em}

\section{Memory Jogger}
In \href{https://www.courtlistener.com/opinion/108838/miller-v-california/}{\textit{Miller v. California}}, the Supreme Court reaffirmed that obscene speech lies outside First Amendment protection and replaced the ``utterly without redeeming social value'' test with a three-prong obscenity framework anchored in contemporary local community standards and lacking serious literary, artistic, political, or scientific value.

\section{Detailed Case Facts}
Marvin Miller, who operated a California mail-order business distributing adult books and films, conducted a mass mailing campaign advertising sexually explicit illustrated books and a film (\href{https://supreme.justia.com/cases/federal/us/413/15/}{Justia, 413 U.S. at 16--18}). His prosecution stemmed specifically from causing five unsolicited advertising brochures to be mailed in an envelope addressed to a restaurant in Newport Beach, California (\href{https://supreme.justia.com/cases/federal/us/413/15/}{413 U.S. at 18}). The restaurant manager and his mother opened the envelope; having never requested the materials, they complained to local police (\href{https://supreme.justia.com/cases/federal/us/413/15/}{413 U.S. at 18}).

The brochures advertised four books---\textit{Intercourse}, \textit{Man-Woman}, \textit{Sex Orgies Illustrated}, and \textit{An Illustrated History of Pornography}---and a film titled \textit{Marital Intercourse} (\href{https://supreme.justia.com/cases/federal/us/413/15/}{413 U.S. at 18}). Although the brochures included some descriptive text, they consisted primarily of photographs and drawings explicitly depicting men and women in groups of two or more engaging in sexual activities, with genitals prominently displayed (\href{https://supreme.justia.com/cases/federal/us/413/15/}{413 U.S. at 18}).

Miller was arrested and charged with a misdemeanor under California Penal Code \S~311.2(a), which prohibited knowingly sending, distributing, or offering to distribute obscene matter (\href{https://supreme.justia.com/cases/federal/us/413/15/}{413 U.S. at 16--17 \& n.1}). At the time of the offense (prior to a June 25, 1969 statutory amendment), California Penal Code \S~311(a) defined ``obscene'' largely along the lines of the plurality test in \href{https://supreme.justia.com/cases/federal/us/383/413/}{\textit{Memoirs v. Massachusetts}}, 383 U.S. 413 (1966), requiring that the predominant appeal of the matter, taken as a whole, be to prurient interest, go substantially beyond customary limits of candor, and be ``utterly without redeeming social importance'' (\href{https://supreme.justia.com/cases/federal/us/413/15/}{413 U.S. at 17 n.1, 23}).

\section{Procedural History}
Miller was tried before a jury in Orange County, California (\href{https://supreme.justia.com/cases/federal/us/413/15/}{413 U.S. at 16--17, 31}). Both the prosecution and the defense tried the case on the assumption that California's statewide standards governed; the trial judge instructed the jury to evaluate prurient appeal and patent offensiveness under the ``contemporary community standards of the State of California'' (\href{https://supreme.justia.com/cases/federal/us/413/15/}{413 U.S. at 31}). Defense counsel did not object at trial to the statewide instruction or to the testimony of the State's police expert on California community standards (\href{https://supreme.justia.com/cases/federal/us/413/15/}{413 U.S. at 31 \& n.12}). The jury convicted Miller of violating California Penal Code \S~311.2(a) (\href{https://supreme.justia.com/cases/federal/us/413/15/}{413 U.S. at 16}).

On appeal to the Appellate Department of the Superior Court of California, County of Orange, Miller contended for the first time that instructing the jury on state rather than national community standards violated the First and Fourteenth Amendments (\href{https://supreme.justia.com/cases/federal/us/413/15/}{413 U.S. at 31}). The Appellate Department summarily affirmed the conviction without a published opinion (\href{https://supreme.justia.com/cases/federal/us/413/15/}{413 U.S. at 16--17}). Miller appealed to the U.S. Supreme Court, which heard oral argument on January 18--19, 1972, ordered reargument on November 7, 1972, and issued its decision on June 21, 1973, alongside companion obscenity cases including \textit{Paris Adult Theatre I v. Slaton}, 413 U.S. 49 (1973) (\href{https://www.oyez.org/cases/1971/70-73}{Oyez}; \href{https://supreme.justia.com/cases/federal/us/413/15/}{413 U.S. at 15, 27}).

\section{Judicial Votes}
\textbf{Vote:} 5--4, vacating and remanding the judgment of the Appellate Department of the Superior Court of California, County of Orange (\href{https://supreme.justia.com/cases/federal/us/413/15/}{413 U.S. at 16, 37}).
\begin{itemize}
\item \textbf{Majority (5):} Chief Justice Warren E. Burger (author), joined by Justices Byron R. White, Harry A. Blackmun, Lewis F. Powell, Jr., and William H. Rehnquist.
\item \textbf{Concurrences (0):} No concurring opinions were filed.
\item \textbf{Dissents (4):}
\begin{itemize}
\item Justice William O. Douglas filed a solo dissenting opinion.
\item Justice William J. Brennan, Jr., filed a dissenting opinion, joined by Justices Potter Stewart and Thurgood Marshall.
\end{itemize}
\end{itemize}

\section{Holding}
\textbf{Judgment:} Vacated and remanded to the Appellate Department of the Superior Court, Orange County, California, for further proceedings consistent with the newly announced First Amendment obscenity standards (\href{https://supreme.justia.com/cases/federal/us/413/15/}{413 U.S. at 37}).

\textbf{Rule Announced:} Reaffirming \textit{Roth v. United States}, 354 U.S. 476 (1957), the Court held that obscene material is categorically unprotected by the First Amendment (applicable to the States through the Fourteenth Amendment) (\href{https://supreme.justia.com/cases/federal/us/413/15/}{413 U.S. at 23, 36}). State regulation of obscenity is constitutionally permissible only when limited to works depicting or describing sexual conduct specifically defined by applicable state law (as written or authoritatively construed) and evaluated under three conjunctive guidelines for the trier of fact (\href{https://supreme.justia.com/cases/federal/us/413/15/}{413 U.S. at 24}):
\begin{enumerate}
\item[(a)] Whether ``the average person, applying contemporary community standards'' would find that the work, taken as a whole, appeals to the prurient interest;
\item[(b)] Whether the work depicts or describes, in a patently offensive way, sexual conduct specifically defined by the applicable state law; and
\item[(c)] Whether the work, taken as a whole, lacks serious literary, artistic, political, or scientific value (commonly abbreviated ``SLAPS'' value).
\end{enumerate}
The Court rejected the \href{https://supreme.justia.com/cases/federal/us/383/413/}{\textit{Memoirs}} plurality's ``utterly without redeeming social value'' requirement and held that juries may measure prurient appeal and patent offensiveness by local or statewide community standards rather than a hypothetical ``national standard'' (\href{https://supreme.justia.com/cases/federal/us/413/15/}{413 U.S. at 24--25, 30--34, 36--37}).

\textbf{Holding vs.\ Dicta:} The core holding encompasses four points: (1) obscenity remains unprotected under the First Amendment; (2) the three-part \textit{Miller} test governs obscenity determinations, superseding \textit{Memoirs}'s ``utterly without redeeming social value'' test; (3) the regulated sexual conduct must be specifically defined by state law as written or authoritatively construed; and (4) the First Amendment permits factfinders to apply contemporary state or local community standards rather than a uniform national standard to prongs (a) and (b) (\href{https://supreme.justia.com/cases/federal/us/413/15/}{413 U.S. at 23--25, 30--34}). By contrast, Chief Justice Burger's two ``plain examples'' of regulable hard-core sexual conduct---(a) patently offensive representations of ultimate sexual acts and (b) patently offensive representations of masturbation, excretory functions, and lewd exhibition of the genitals---are illustrative dicta explicitly framed as non-exhaustive guidance rather than a legislative code (\href{https://supreme.justia.com/cases/federal/us/413/15/}{413 U.S. at 25}). Likewise, his historical remarks asserting that nineteenth-century censorship did not chill American belles lettres and his observation in Footnote 8 that States enjoy broader power over live nonverbal conduct under \textit{United States v. O'Brien}, 391 U.S. 367 (1968), are dicta (\href{https://supreme.justia.com/cases/federal/us/413/15/}{413 U.S. at 26 n.8, 35--36}).

\section{Analysis of Opinions}

\subsection{Majority Opinion (Chief Justice Burger)}
Chief Justice Burger framed the case as an opportunity to bring majority consensus to ``the intractable obscenity problem'' after years of fractured post-\textit{Roth} rulings and summary per curiam reversals under \textit{Redrup v. New York}, 386 U.S. 767 (1967), which had cast the Supreme Court as ``an unreviewable board of censorship for the 50 States'' (\href{https://supreme.justia.com/cases/federal/us/413/15/}{413 U.S. at 16, 22 \& n.3, 29}). He emphasized that Miller's brochures were thrust upon unwilling recipients via aggressive mass mailing, implicating the States' recognized interest in shielding unconsenting adults and juveniles from obscene material (\href{https://supreme.justia.com/cases/federal/us/413/15/}{413 U.S. at 18--19}).

Burger then dismantled the \href{https://supreme.justia.com/cases/federal/us/383/413/}{\textit{Memoirs}} plurality test. Whereas \textit{Roth} had treated obscenity as presumptively without redeeming social importance, the three-Justice \textit{Memoirs} plurality forced prosecutors to prove a negative beyond a reasonable doubt---that the work was ``utterly without redeeming social value''---an impossible burden that never commanded a Court majority and had now been abandoned even by its author, Justice Brennan (\href{https://supreme.justia.com/cases/federal/us/413/15/}{413 U.S. at 21--25}). To isolate ``hard core'' pornography while safeguarding serious expression and ensuring fair notice, Burger announced the three-prong \textit{Miller} test, coupling the requirement that state law specifically define the proscribed sexual conduct with the guarantee of independent appellate review of constitutional claims (\href{https://supreme.justia.com/cases/federal/us/413/15/}{413 U.S. at 24--27}).

On community standards, Burger rejected any constitutional requirement of a single uniform ``national standard'' for prongs (a) and (b). Because prurient appeal and patent offensiveness are essentially factual questions in a vast, culturally diverse nation, requiring juries to divine an abstract national standard would be ``an exercise in futility'' (\href{https://supreme.justia.com/cases/federal/us/413/15/}{413 U.S. at 30}). Relying on Chief Justice Warren's \textit{Jacobellis v. Ohio} dissent and Justice Harlan's federalism concerns, Burger concluded that the First Amendment does not force citizens in Maine or Mississippi to tolerate public depictions accepted in Las Vegas or New York City, and that instructing the jury on California statewide standards adequately protected against judging material by the sensibilities of an unusually susceptible person (\href{https://supreme.justia.com/cases/federal/us/413/15/}{413 U.S. at 32--34}). Finally, Burger dismissed the dissents' warnings of repression, distinguishing the ``unfettered interchange of ideas'' from ``commercial exploitation of obscene material'' (\href{https://supreme.justia.com/cases/federal/us/413/15/}{413 U.S. at 34--36}).

\subsection{Dissenting Opinion (Justice Douglas)}
Justice Douglas mounted a sweeping textualist and due process attack on both the Court's new test and the entire obscenity exception. First, he argued that ``obscenity'' appears nowhere in the Constitution or the Bill of Rights and that no historical exception at the Founding excluded allegedly obscene publications from the First Amendment's protection of ``the press'' (\href{https://supreme.justia.com/cases/federal/us/413/15/}{413 U.S. at 40}). Without textual benchmarks, judges merely enforce personal ``tastes and standards of literature''; if society wishes to authorize censorship, it must do so by constitutional amendment rather than judicial improvisation (\href{https://supreme.justia.com/cases/federal/us/413/15/}{413 U.S. at 40--41, 46--47}).

Second, Douglas stressed fair warning and vagueness. Applying freshly minted constitutional standards to uphold a criminal prosecution retroactively carried ``all the evils of an \textit{ex post facto} law'' (\href{https://supreme.justia.com/cases/federal/us/413/15/}{413 U.S. at 37, 41--42}). At a minimum, Douglas contended, due process barred criminal punishment unless a specific work had first been adjudicated obscene in a completed civil declaratory proceeding (\href{https://supreme.justia.com/cases/federal/us/413/15/}{413 U.S. at 41--44}). Third, he argued that using ``offensive'' as a criminal standard cut ``the very vitals out of the First Amendment,'' noting that the Court had struck down a municipal ordinance punishing ``annoying'' assembly as unconstitutionally vague in \textit{Coates v. City of Cincinnati}, 402 U.S. 611 (1971) (\href{https://supreme.justia.com/cases/federal/us/413/15/}{413 U.S. at 44--46}).

\subsection{Dissenting Opinion (Justice Brennan, joined by Justices Stewart and Marshall)}
Justice Brennan---the original author of \textit{Roth} and the \textit{Memoirs} plurality---incorporated his extensive dissent in the companion case \textit{Paris Adult Theatre I v. Slaton}, 413 U.S. 49, 73 (1973), where he formally abandoned the effort to define obscenity for consenting adults because no judicial formulation could separate obscene from protected expression with sufficient precision to avoid vagueness, chilling effects, and institutional stress on the courts (\href{https://supreme.justia.com/cases/federal/us/413/15/}{413 U.S. at 27--29, 47}; \href{https://caselaw.findlaw.com/court/us-supreme-court/425/963.html}{\textit{Liles v. Oregon}, 425 U.S. 963 (1976)}).

Although Brennan reserved judgment on whether a narrowly drawn statute could constitutionally punish distribution to juveniles or obtrusive exposure to unconsenting adults (such as Miller's unsolicited mailing), he concluded that California Penal Code \S~311.2(a) was facially invalid for substantial overbreadth because it swept in all distribution of obscene matter, including consensual adult transactions (\href{https://supreme.justia.com/cases/federal/us/413/15/}{413 U.S. at 47--48}). Invoking First Amendment overbreadth doctrine (\textit{Gooding v. Wilson} and \textit{Dombrowski v. Pfister}), Brennan argued that Miller did not need to show his own conduct was protected to challenge a facially overbroad statute, and he would have reversed and remanded for the state courts to determine whether a narrowing construction could rehabilitate the law (\href{https://supreme.justia.com/cases/federal/us/413/15/}{413 U.S. at 47--48}).

\section{Examples: Future Applications}
\begin{enumerate}
\item \textbf{Same-Side (Unprotected Obscenity):} \textit{A commercial distributor mails unsolicited, graphically explicit pamphlets depicting actual sexual intercourse and masturbation, with no narrative or educational commentary, to random households in a conservative county whose state obscenity statute specifically enumerates those acts.} This falls squarely on the \textit{Miller} side: under local community standards, the pamphlets appeal to a prurient interest, depict statutorily defined hard-core sexual conduct in a patently offensive way, and lack serious literary, artistic, political, or scientific value (\href{https://supreme.justia.com/cases/federal/us/413/15/}{413 U.S. at 24--25}).

\item \textbf{Same-Side (Unprotected Obscenity):} \textit{A publisher inserts a single quotation from Voltaire on the opening page of an otherwise plotless, hard-core pornographic magazine consisting solely of explicit photographs of sexual acts and lewd genital displays.} Under \textit{Miller}, the work is judged ``taken as a whole'' and must possess \textit{serious} SLAPS value; Burger explicitly noted in Footnote 7 that a token philosophical quotation ``will not constitutionally redeem an otherwise obscene publication,'' rejecting \textit{Memoirs}'s ``utterly without'' floor (\href{https://supreme.justia.com/cases/federal/us/413/15/}{413 U.S. at 24--25 \& n.7}).

\item \textbf{Opposite-Side (Protected Speech):} \textit{A state prosecutor charges a medical textbook publisher under an obscenity statute because an obstetrics and human-sexuality manual contains graphic anatomical photographs and explicit descriptions of sexual organs and reproductive acts.} Even if a local jury found the imagery offensive, the textbook is constitutionally protected under \textit{Miller} prong (c) because, taken as a whole, it possesses serious scientific and educational value---an example Burger expressly invoked as protected (\href{https://supreme.justia.com/cases/federal/us/413/15/}{413 U.S. at 26}; see also \href{https://supreme.justia.com/cases/federal/us/481/497/}{\textit{Pope v. Illinois}, 481 U.S. 497 (1987)}, holding prong (c) is judged by an objective reasonable-person standard, not local community standards).

\item \textbf{Opposite-Side (Protected Speech):} \textit{A state enforces a broadly worded statute criminalizing ``lewd, indecent, or immoral performances'' against a critically acclaimed theatrical film containing non-explicit nudity and simulated romance, where the statute neither defines specific sexual conduct nor has been authoritatively narrowed by state courts.} This fails \textit{Miller} on multiple grounds: prong (b) requires patently offensive depictions of hard-core sexual conduct \textit{specifically defined} by state law, mere nudity is not obscene, and serious artistic value under prong (c) independently shields the film (\href{https://supreme.justia.com/cases/federal/us/413/15/}{413 U.S. at 24--26}; \href{https://supreme.justia.com/cases/federal/us/535/564/}{\textit{Ashcroft v. ACLU}, 535 U.S. 564, 576--79 (2002)}).

\item \textbf{Fence-Sitter (Close Case):} \textit{An independent avant-garde filmmaker streams a sexually explicit online film that includes unsimulated sexual acts alongside an extended feminist political monologue critiquing consumer culture, and federal prosecutors bring obscenity charges in a rural district where the stream was downloaded.} This is a classic fence-sitter: prongs (a) and (b) may be satisfied under the forum's local community standards, raising the unresolved Internet-era tension over whether online speakers are subjected to the most conservative jurisdiction's norms (\href{https://supreme.justia.com/cases/federal/us/535/564/}{\textit{Ashcroft v. ACLU}, 535 U.S. at 580--85}; \href{https://firstamendment.mtsu.edu/article/community-standards/}{First Amendment Encyclopedia}). At the same time, prong (c) turns on whether a reasonable person would find the political monologue confers \textit{serious} artistic or political value on the work ``taken as a whole'' rather than serving as a pretextual wrapper (\href{https://supreme.justia.com/cases/federal/us/481/497/}{\textit{Pope}, 481 U.S. at 500--01}).
\end{enumerate}

\section{Critique}
\textbf{Originalist and Textualist Perspectives:} Originalists and textualists divide sharply over \textit{Miller}. Defenders of the categorical obscenity exception emphasize that the First Amendment was never understood as absolute and point to Founding-era and nineteenth-century common-law and statutory restrictions on obscene publications as historical evidence that obscenity lay outside ``the freedom of speech'' (\href{https://supreme.justia.com/cases/federal/us/413/15/}{413 U.S. at 23, 29, 35}; \href{https://www.fire.org/research-learn/obscenity-exception-first-amendment}{FIRE}). They also applaud \textit{Miller}'s federalism: decentralizing prongs (a) and (b) to local juries respects state police powers and avoids judicial imposition of a single national moral code (\href{https://supreme.justia.com/cases/federal/us/413/15/}{413 U.S. at 30--33}). Conversely, strict textualists echoing Justice Douglas argue that the First Amendment's text contains no obscenity carve-out, that Founding-era obscenity prosecutions were rare and largely tied to blasphemy or seditious libel, and that \textit{Miller}'s multi-prong balancing test is an unmoored judicial invention that invites subjective censorship (\href{https://supreme.justia.com/cases/federal/us/413/15/}{413 U.S. at 40--41}; \href{https://www.law.cornell.edu/constitution-conan/amendment-1/overview-of-obscene-speech}{Cornell LII, Constitution Annotated}). Even within the conservative wing, Justice Scalia later questioned whether ``reasonable person'' evaluation of artistic value under \href{https://supreme.justia.com/cases/federal/us/481/497/}{\textit{Pope v. Illinois}} is coherent, remarking that ``just as there is no use arguing about taste, there is no use litigating about it.''

\textbf{Progressive, Feminist, and Doctrinal Perspectives:} Civil-libertarian and doctrinal critics contend that \textit{Miller} suffers from chronic vagueness and produces a severe chilling effect on publishers, librarians, and artists, who cannot predict how distant juries will apply ``prurient interest'' or ``patently offensive'' (\href{https://firstamendment.mtsu.edu/article/miller-v-california/}{First Amendment Encyclopedia}; \href{https://journals.library.wustl.edu/jurisprudence/article/8891/galley/25674/view/}{Wash.\ U.\ Jurisprudence Rev.}). The rise of the Internet magnified this flaw: because online publishers cannot easily geo-fence content, applying local ``contemporary community standards'' risks subjecting national and global speech to the lowest common denominator of the most restrictive forum, an issue that fractured the Court in \href{https://supreme.justia.com/cases/federal/us/535/564/}{\textit{Ashcroft v. ACLU}, 535 U.S. 564 (2002)} (\href{https://firstamendment.mtsu.edu/article/community-standards/}{First Amendment Encyclopedia}). From a different progressive angle, feminist legal theorist Catharine MacKinnon criticized \textit{Miller} from the opposite direction, arguing that obscenity law focuses on moral offense and prudery rather than the real-world civil-rights harm of pornography---namely, the sexual subordination and inequality of women (\href{https://firstamendment.mtsu.edu/article/catharine-mackinnon/}{First Amendment Encyclopedia}; \href{https://virginialawreview.org/wp-content/uploads/2020/12/MacKinnon_Book.pdf}{Va.\ L.\ Rev.}).

\textbf{Underlying Competing Values:} \textit{Miller} sits at the intersection of four enduring constitutional tensions: (1) individual expressive autonomy and protection against majoritarian censorship versus community self-determination and moral order; (2) uniform national First Amendment baselines versus federalism and local cultural diversity; (3) fair notice and anti-vagueness in criminal law versus the practical impossibility of drafting exhaustive statutory definitions of sexual depictions; and (4) the boundary between high-value political or artistic discourse and commercialized, non-cognitive sexual exploitation.

\section{Key Quotations}
\begin{quote}
``The basic guidelines for the trier of fact must be: (a) whether `the average person, applying contemporary community standards' would find that the work, taken as a whole, appeals to the prurient interest, \textit{Kois v. Wisconsin}, supra, at 408 U. S. 230, quoting \textit{Roth v. United States}, supra, at 354 U. S. 489; (b) whether the work depicts or describes, in a patently offensive way, sexual conduct specifically defined by the applicable state law; and (c) whether the work, taken as a whole, lacks serious literary, artistic, political, or scientific value.''
\end{quote}
\noindent --- Chief Justice Warren E. Burger, Opinion of the Court, 413 U.S. at 24 (\href{https://supreme.justia.com/cases/federal/us/413/15/}{Justia}; \href{https://www.courtlistener.com/opinion/108838/miller-v-california/}{CourtListener}).

\begin{quote}
``It is neither realistic nor constitutionally sound to read the First Amendment as requiring that the people of Maine or Mississippi accept public depiction of conduct found tolerable in Las Vegas, or New York City.''
\end{quote}
\noindent --- Chief Justice Warren E. Burger, Opinion of the Court, 413 U.S. at 32 (\href{https://supreme.justia.com/cases/federal/us/413/15/}{Justia}; \href{https://www.courtlistener.com/opinion/108838/miller-v-california/}{CourtListener}).

\begin{quote}
``The First Amendment protects works which, taken as a whole, have serious literary, artistic, political, or scientific value, regardless of whether the government or a majority of the people approve of the ideas these works represent.''
\end{quote}
\noindent --- Chief Justice Warren E. Burger, Opinion of the Court, 413 U.S. at 34 (\href{https://supreme.justia.com/cases/federal/us/413/15/}{Justia}; \href{https://www.courtlistener.com/opinion/108838/miller-v-california/}{CourtListener}).

\begin{quote}
``The difficulty is that we do not deal with constitutional terms, since `obscenity' is not mentioned in the Constitution or Bill of Rights. And the First Amendment makes no such exception from `the press' which it undertakes to protect nor, as I have said on other occasions, is an exception necessarily implied, for there was no recognized exception to the free press at the time the Bill of Rights was adopted which treated `obscene' publications differently from other types of papers, magazines, and books. So there are no constitutional guidelines for deciding what is and what is not `obscene.'\,''
\end{quote}
\noindent --- Justice William O. Douglas, Dissenting, 413 U.S. at 40 (\href{https://supreme.justia.com/cases/federal/us/413/15/}{Justia}; \href{https://www.courtlistener.com/opinion/108838/miller-v-california/}{CourtListener}).

\begin{quote}
``The idea that the First Amendment permits punishment for ideas that are `offensive' to the particular judge or jury sitting in judgment is astounding. No greater leveler of speech or literature has ever been designed. To give the power to the censor, as we do today, is to make a sharp and radical break with the traditions of a free society. The First Amendment was not fashioned as a vehicle for dispensing tranquilizers to the people.''
\end{quote}
\noindent --- Justice William O. Douglas, Dissenting, 413 U.S. at 44--45 (\href{https://supreme.justia.com/cases/federal/us/413/15/}{Justia}; \href{https://www.courtlistener.com/opinion/108838/miller-v-california/}{CourtListener}).

\end{document}
